\section{Related Work}
\subsection{LoRa and Mesh Networking}
LoRa is a low-power wide-area technology whose range can reach several kilometers depending on the antenna, radio configuration, terrain, and environment \cite{centenaro2016longrange}, which makes it a good fit for low-bandwidth, low-cost, scalable monitoring \cite{rahman2018lora}. Multi-hop and mesh operation can extend coverage further: Andersen et al.\ approximately doubled the range of an asset-tracking prototype with a LoRa mesh \cite{andersen2020mesh}, and Mullick et al.\ reported improved air-quality monitoring under adverse conditions \cite{mullick2022mesh}. Additional hops can increase latency because nodes must coordinate transmissions \cite{centelles2021multihop}, but relays shorten the distance of each transmission and may improve energy efficiency \cite{khonrang2022multihop}.

\subsection{Low-Power Design}
For battery-powered nodes with low duty cycles, the sleep period usually matters more for energy use than the active period \cite{hwang2012poweraware}. Kombo et al.\ reduced the duty cycle of a monitoring system to below one percent \cite{kombo2021groundwater}, and Hung et al.\ reported power reductions of about 38.46--54.44\% with a low-power wireless sensor protocol \cite{hung2020lowpower}. This motivates scheduled sensing, switching off peripherals, and sleeping the radio, as used here.

\subsection{Soil Moisture Sensing}
Gravimetric measurement, which compares the weight of a soil sample before and after oven drying, is simple and accurate and is the usual calibration reference, but it is destructive and unsuitable for continuous monitoring \cite{ref24,ref25}. TDR and FDR are accurate and non-destructive but expensive \cite{ref19,ref28}. Resistive sensors are cheap but unreliable across soil types and prone to corrosion \cite{ref28}. Capacitive sensors respond to the dielectric constant of the soil, are inexpensive, easy to install, and resistant to corrosion, and are therefore widely preferred for low-cost deployments \cite{ref33,ref34,ref35,ref36}. They need calibration for the soil in use, which this work performs against gravimetric reference data.

\subsection{Cloud Platforms and Gap}
Cloud platforms let researchers access time-series data from remote sites without retrieving the devices \cite{naik2017cloud,sheffield2022cloud}. The individual technologies above are well established. This work does not propose a new sensing principle or radio technology. It combines them into an affordable, buried, unattended monitoring platform for turf roofs, with local backup and cloud access, and evaluates it in the field.
